\usepackage[LGR,T1]{fontenc}
\usepackage[toc,page]{appendix} 	% appendices
\usepackage{mathptmx}
\usepackage{mathtools,amssymb,amsthm,amsfonts,amsmath}
\usepackage{enumerate} 	
\usepackage{float}
\usepackage[shortlabels]{enumitem}			
\usepackage[hidelinks]{hyperref} 	
\usepackage{bussproofs}				% proof trees
\usepackage{longtable}
\usepackage{multicol}
\usepackage{authblk} % for affiliation
\usepackage{enumerate} % roman enumeration
\usepackage[hidelinks]{hyperref} % hyperref without ugly box borders
\usepackage{url}
\usepackage{cite}
\usepackage{cleveref} % for cross-referencing
\usepackage{textgreek}				% greek letters
\usepackage{chngcntr}				% continuous section numbers

\allowdisplaybreaks

\newcommand{\red}[1]{{\color{red} #1}}

\newcommand{\textgreek}[1]{\begingroup\fontencoding{LGR}\selectfont#1\endgroup}

\DeclareMathAlphabet{\mathcal}{OMS}{cmsy}{m}{n} % for regular mathcal font


%\usepackage[
%	paperwidth=6.10in,
%	paperheight=9.25in,
%	left=1.0in,
%	right=1.0in,
%	top=1.0in,
%	bottom=1.0in,
%]{geometry}


\newtheorem{theorem}{Theorem}[section]
\newtheorem{lemma}[theorem]{Lemma}
\newtheorem{corollary}[theorem]{Corollary}
\newtheorem{definition}[theorem]{Definition}
\newtheorem{example}[theorem]{Example}
\newtheorem{remark}{Remark}
\newtheorem{exercise}[theorem]{Exercise}


%
%
%\usepackage{parskip}
%\setlength{\parskip}{\baselineskip}
%%\setlength{\parskip}{2em}
%\linespread{1.20} 

\usepackage{stmaryrd} % for llbrackets and rrbrackets

\usepackage[table,dvipsnames]{xcolor}

\usepackage{tikz-cd}
\usetikzlibrary{patterns}
\tikzcdset{every label/.append style = {font = \small}} %changes label size
\usetikzlibrary{decorations.markings}

\usepackage{setspace}				% for spacing

% my commands

\renewcommand{\phi}{\varphi}

\newcommand{\J}{\mathcal{J}}

\newcommand{\gga}{\textit{Grundgesetze}}
\newcommand{\bs}{{\it Begriffsschrift}}
\newcommand{\gla}{\textit{Grundlagen}}
\newcommand{\wiss}{\textit{Wissenschaftslehre}}
\newcommand{\zfc}{\textsf{ZF}}

\newcommand{\LU}{\textit{Logische Untersuchungen}}
\newcommand{\FuB}{\textit{Funktion und Begriff}}



\newcommand{\prop}{\> \mathsf{prop}}
\newcommand{\type}{\mathtt{type}}
\newcommand{\true}{\> \mathsf{true}}
\newcommand{\val}{\> \mathsf{val}}
\newcommand{\ctx}{\> \mathsf{ctx}}
\newcommand{\pitype}[1]{\Pi_{(#1)}}
\newcommand{\funtype}[2]{#1\to#2}
\newcommand{\infuntype}[2]{(#1\to#2)}

\newcommand{\hypjdg}[2]{#1 \gg #2} %{#2 \; (#1)}

\newcommand{\fun}[1]{\lambda #1.}

\newcommand{\app}{\mathsf{app}}
\newcommand{\void}{\mathtt{void}}
\newcommand{\unit}{\mathtt{unit}}
\renewcommand{\succ}{\mathsf{s}}

\newcommand{\expr}{\mathsf{expr}}
\newcommand{\var}{\mathsf{var}}

\newcommand{\bool}{\mathtt{bool}}


\newcommand{\nat}{\mathsf{nat}}%{\mathbb{N}}
\newcommand{\zero}{\mathsf{0}}
\newcommand{\W}{\mathsf{W}} %\mathsf{nat}
%\newcommand{\natrec}{\mathsf{natrec}}


%\newcommand{\refl}{\mathsf{refl}}
\newcommand{\eqrec}{\mathsf{eqrec}}
\newcommand{\shouldsf}[1]{#1}

\renewcommand{\tt}{\mathsf{true}}
\newcommand{\ff}{\mathsf{false}}
\newcommand{\ite}[4]{\mathsf{ite}(#1, #2, #3)} 


\newcommand{\plus}{\mathsf{+}}

\newcommand{\pair}[2]{\mathsf{pair}({#1,#2})}%{\langle #1,#2 \rangle}
\newcommand{\fst}[1]{\textsf{fst}(#1)}
\newcommand{\snd}[1]{\textsf{snd}(#1)}
\newcommand{\spread}[2]{\textsf{spread}(#1, #2)}
\newcommand{\inl}[1]{\textsf{inl}(#1)}
\newcommand{\inr}[1]{\textsf{inr}(#1)}
\newcommand{\case}[3]{\textsf{case}(#1, #2, #3)}
\newcommand{\abort}[1]{\textsf{abort}(#1)}
%\newcommand{\spread}[2]{\textsf{spread}(#1, #2)}

\newcommand{\U}{\mathcal{U}}
\newcommand{\urec}[4]{\mathsf{urec}(#1, #2, #3)} 

\newcommand{\wrec}[4]{\mathsf{wrec}(#1, #2, #3)} 
\renewcommand{\sup}[2]{\mathsf{sup}(#1, #2)}

\renewcommand{\succ}[1]{\mathsf{s}(#1)}
%\renewcommand{\succop}{\mathsf{s}}
\newcommand{\natrec}[4]{\mathsf{rec}(#1, #2, #3)} 
\newcommand{\natrecop}{\mathsf{rec}} 
%\newcommand{\reflop}[1]{\mathsf{refl}(#1)}
\newcommand{\idrec}[5]{\mathsf{J}(#1, #2, #3, #4)} %{\mathsf{idrec}_{#5}(#1, #2, #3, #4)}
\newcommand{\rfl}{\mathtt{rfl}}
\newcommand{\subs}[2]{[#1/#2]}
\newcommand{\idtype}[3]{#2 = {#1}\ #3}

\newcommand{\univ}[1]{\mathcal{#1}}
\newcommand*\sq{\mathbin{\vcenter{\hbox{\rule{.3ex}{.3ex}}}}}
\newcommand{\ua}{\mathsf{ua}}
\newcommand{\idtoeqv}{\mathsf{idtoeqv}}

\newcommand{\ucircle}{\mathsf{S}^1}
\newcommand{\base}{\mathsf{base}}
\newcommand{\uloop}{\mathsf{loop}}

%% Hyphenation

\tolerance=5000
\emergencystretch=\maxdimen
\hyphenpenalty=5000 
\hbadness=10000

% Footnote

\renewcommand*\footnoterule{}
\setlength{\skip\footins}{1.5pc plus 5pt minus 2pt}

% Notations

\DeclareFontFamily{U}{matha}{\hyphenchar\font45}
\DeclareFontShape{U}{matha}{m}{n}{
	<5> <6> <7> <8> <9> <10> gen * matha
	<10.95> matha10 <12> <14.4> <17.28> <20.74> <24.88> matha12
}{}
\DeclareSymbolFont{matha}{U}{matha}{m}{n}

\DeclareMathSymbol{\supvar}{3}{matha}{"81} % inspired by https://tex.stackexchange.com/questions/14386/importing-a-single-symbol-from-a-different-font

\makeatletter
\newcommand{\citecomment}[2][]{\citen{#2}#1\citevar}
\newcommand{\citeone}[1]{\citecomment{#1}}
\newcommand{\citetwo}[2][]{\citecomment[,~#1]{#2}}
\newcommand{\citevar}{\@ifnextchar\bgroup{;~\citeone}{\@ifnextchar[{;~\citetwo}{]}}}
\newcommand{\citefirst}{\@ifnextchar\bgroup{\citeone}{\@ifnextchar[{\citetwo}{]}}}
\newcommand{\cites}{[\citefirst}
\makeatother

\newcommand{\impl}{\supvar}

\DeclareMathSymbol{\mhyphen}{\mathord}{AMSa}{"39}

\usepackage{xcolor,listings}
\usepackage{doc}

\definecolor{keywordcolor}{rgb}{0.5, 0.3, 0.7}   % amethyst
\definecolor{tacticcolor}{rgb}{0.0, 0.1, 0.6}    % blue
\definecolor{commentcolor}{rgb}{0.4, 0.4, 0.4}   % grey
\definecolor{symbolcolor}{rgb}{1.2, 0.4, 0.1}    % orange
\definecolor{sortcolor}{rgb}{0.1, 0.5, 0.1}      % green
\definecolor{attributecolor}{rgb}{0.7, 0.1, 0.1} % red

\newcommand{\simnot}{\mathord{\sim}}

\definecolor{darkgreen}{rgb}{0.18,0.54,0.34}
\definecolor{darkpink}{rgb}{0.75,0.25,0.5}

\makeatletter
\lst@InputCatcodes
\def\lst@DefEC{%
	\lst@CCECUse \lst@ProcessLetter
	^^80^^81^^82^^83^^84^^85^^86^^87^^88^^89^^8a^^8b^^8c^^8d^^8e^^8f%
	^^90^^91^^92^^93^^94^^95^^96^^97^^98^^99^^9a^^9b^^9c^^9d^^9e^^9f%
	^^a0^^a1^^a2^^a3^^a4^^a5^^a6^^a7^^a8^^a9^^aa^^ab^^ac^^ad^^ae^^af%
	^^b0^^b1^^b2^^b3^^b4^^b5^^b6^^b7^^b8^^b9^^ba^^bb^^bc^^bd^^be^^bf%
	^^c0^^c1^^c2^^c3^^c4^^c5^^c6^^c7^^c8^^c9^^ca^^cb^^cc^^cd^^ce^^cf%
	^^d0^^d1^^d2^^d3^^d4^^d5^^d6^^d7^^d8^^d9^^da^^db^^dc^^dd^^de^^df%
	^^e0^^e1^^e2^^e3^^e4^^e5^^e6^^e7^^e8^^e9^^ea^^eb^^ec^^ed^^ee^^ef%
	^^f0^^f1^^f2^^f3^^f4^^f5^^f6^^f7^^f8^^f9^^fa^^fb^^fc^^fd^^fe^^ff%
	^^^^0393^^^^2081^^^^2082^^^^2208^^^^03c3%
	^^^^2113^^^^2192^^^^03a0^^^^03bb%          <-- NEW
	^^00}
\lst@RestoreCatcodes
\makeatother

\lstdefinestyle{pearl}{
	inputencoding=utf8,
	extendedchars=true,
	belowcaptionskip=1\baselineskip,
	breaklines=true,
	mathescape=true,
	keywordstyle=[1]{\ttfamily\color{keywordcolor}},
	keywordstyle=[2]{\ttfamily\color{sortcolor}},
	keywordstyle=[3]{\ttfamily\color{tacticcolor}},
	keywordstyle=[4]{\ttfamily\color{attributecolor}},
	commentstyle=\color{blue}\textit,
	stringstyle=\color{pink}\ttfamily,
	basicstyle=\footnotesize\ttfamily,
	showstringspaces=false,
	identifierstyle={\color{black}}, 
	%
	literate=*
	{|}{{{\color{darkgreen}$|$}}}{1}
	{¬}{$\neg$}{1}
	{ℓ}{{{\color{darkgreen}$\ell$}}}{1}
	{→}{{{\color{darkgreen}$\to$}}}{1}
	{Π}{{{\color{black}$\Pi$}}}{1}
	{λ}{{{\color{black}$\lambda$}}}{1}
	{:}{{{\color{darkgreen}:}}}{1}
	{=}{{{\color{darkgreen}=\;}}}{1}
	{[}{{{\color{darkgreen}[}}}{1}
	{]}{{{\color{darkgreen}]}}}{1}
	{sigma}{{{\color{black}$\sigma$}}}{1}
	{.7}{{{\color{red}.7}}}{2},%
	%
	morekeywords={prefix,inductive,reflect,import,eval,infer,universe,print,definition,def,lemma,lem,thm,theorem,type,attribute}
}